\section{Data}
\label{sec:data}

\subsection{Asset Universe}
\label{subsec:data_assets}

The equity universe consists of ten large-cap Brazilian stocks listed on the B3
exchange: PETR4, VALE3, ITUB4, BBDC4, ABEV3, WEGE3, RENT3, SUZB3, B3SA3, and BBAS3.
These ten tickers were selected based on three criteria: (i)~continuous listing
throughout the full sample period January~2015 to December~2024, (ii)~membership in
the Ibovespa index for at least eight of the ten years, and (iii)~average daily
traded volume exceeding R\$~100 million, ensuring that the portfolio weights implied
by the optimization are implementable without significant market-impact costs.
The ten assets span six sectors---energy and materials (PETR4, VALE3, SUZB3),
financials (ITUB4, BBDC4, BBAS3, B3SA3), consumer staples (ABEV3), industrials
(WEGE3), and services (RENT3)---providing cross-sector diversification within the
concentrated structure of the Brazilian equity market.

Daily adjusted closing prices were obtained from Yahoo Finance via the
\texttt{yfinance} library.
Log returns were computed as $r_{i,t} = \log(P_{i,t}/P_{i,t-1})$ for each asset~$i$
and trading day~$t$.
The full sample spans January~2015 to December~2024, comprising 2,487 trading days.
Table~\ref{tab:descriptive_stats} reports summary statistics for log returns over
the full period.
The cross-sectional dispersion in annualized returns is substantial: WEGE3 delivers
the highest annualized return at 23.9\%, followed by PETR4 at 26.0\%, while ABEV3
records the lowest at 1.2\%, reflecting both sector dynamics and firm-specific factors.
The equal-weight portfolio achieves a mean daily return of $0.000545$ (annualized:
$\approx13.7\%$) with a daily standard deviation of $0.015521$ and a
CVaR$_{0.95}$ of $-0.0249$.

\subsection{Macroeconomic Variables}
\label{subsec:data_macro}

Three macroeconomic time series are incorporated to characterize the Brazilian
monetary and credit environment.

\textbf{SELIC rate.} The daily SELIC overnight rate (Series 11, BCB-SGS) was
retrieved from the Banco Central do Brasil public API
(\texttt{api.bcb.gov.br}).
The SELIC is the primary monetary-policy instrument in Brazil and its level and
direction are key determinants of equity valuations and credit conditions.

\textbf{Credit spread.} The spread between Brazilian corporate bond yields and
government securities was constructed as a proxy for domestic credit risk.
This variable captures periods of financial stress that are not fully reflected
in interest-rate levels alone.

\textbf{CDS proxy.} A proxy for sovereign credit risk was constructed using
the five-year Brazil CDS mid-quote, normalized to have zero mean and unit
variance over rolling 252-day windows to ensure stationarity in the state vector.

In addition, the Brazilian implied volatility index (IVOL-BR), when available,
is included as a measure of market-wide uncertainty.
All macroeconomic variables are aligned to the equity trading calendar using
forward-fill for days on which the macroeconomic series do not report (e.g.,
holidays on which the bond market is closed but equities trade).

\subsection{SELIC Regime Classification}
\label{subsec:data_regimes}

Monetary-policy regimes are identified from the direction of the SELIC rate
using a rule-based classification.
A regime is labeled \emph{Rising} if the SELIC rate increased in the most recent
Monetary Policy Committee (COPOM) meeting, \emph{Falling} if it decreased, and
\emph{Stable} if it remained unchanged.
This classification is strictly backward-looking---only information available at
each date is used---so it introduces no look-ahead bias.
Table~\ref{tab:regime_stats} summarizes the regime distribution across the
2,487-day sample.

Over the full period, the Stable regime accounts for the majority of trading days
(42.1\%, or 1,048 days), reflecting the extended periods between rate decisions.
The Rising regime covers 23.5\% of the sample (584 days), concentrated primarily
in 2015--2016 and 2021--2023.
The Falling regime accounts for the remaining 34.4\% (855 days), corresponding
to the easing cycle of 2017--2020 and the initial phase of the 2023--2024
normalization.

The economic significance of these regimes is confirmed by the large variation in
aggregate Sharpe ratios: the equal-weight portfolio achieves a Sharpe ratio of
$-0.23$ during Rising-rate periods, $+0.88$ during Stable periods, and $+0.62$
during Falling periods.
This three-way split provides a natural test of whether a regime-aware portfolio
strategy can exploit this heterogeneity to improve risk-adjusted performance.

\subsection{Stationarity}
\label{subsec:data_stationarity}

All ten log-return series were tested for stationarity using the Augmented
Dickey-Fuller (ADF) test.
All ten tickers reject the unit-root null hypothesis at the 1\% significance level
($p < 0.001$), confirming that log returns are stationary in levels.
This is consistent with standard results in the empirical finance literature and
validates the use of log returns as the primary input to the model without
additional differencing.

\inputtable{descriptive_stats}

\begin{figure}[htbp]
  \centering
  \includegraphics[width=\linewidth]{fig_selic_regimes}
  \caption{SELIC overnight rate and identified monetary-policy regimes, 2015--2024.
           Shaded bands indicate Rising (red), Stable (gray), and Falling (blue) rate
           regimes classified from COPOM decisions.}
  \label{fig:selic_regimes}
\end{figure}
